% ResearchBench — ICLR-format paper draft
% Build: pdflatex researchbench_iclr.tex (twice)
% If you obtain the official iclr2026_conference.sty, place it next to this file;
% it will be used automatically. Otherwise a faithful emulation is applied.
\documentclass{article}

\IfFileExists{iclr2026_conference.sty}{%
  \usepackage{iclr2026_conference}%
}{%
  % ---- ICLR-style emulation (matches official geometry & typography) ----
  \usepackage[letterpaper]{geometry}
  \geometry{textheight=9in, textwidth=5.5in, top=1in, headheight=12pt,
            headsep=25pt, footskip=30pt, centering}
  \usepackage{times}
  \usepackage{fancyhdr}
  \pagestyle{fancy} 
  \fancyhf{}
  \fancyhead[L]{\small Under review as a conference paper at ICLR 2027}
  \fancyfoot[C]{\thepage}
  \renewcommand{\headrulewidth}{0.4pt}
  \usepackage{titlesec}
  \titleformat{\section}{\large\bfseries\scshape}{\thesection}{1em}{}
  \titleformat{\subsection}{\normalsize\bfseries\scshape}{\thesubsection}{1em}{}
  \titleformat{\subsubsection}{\normalsize\bfseries}{\thesubsubsection}{1em}{}
}

\usepackage[round,authoryear]{natbib}
\usepackage{booktabs}
\usepackage{hyperref}
\usepackage{graphicx}
\usepackage{amsmath}
\usepackage{url}
\usepackage{microtype}
\usepackage{xcolor}
\newcommand{\todo}[1]{\textcolor{red}{[TODO: #1]}}
\newcommand{\site}{\url{https://infinity-ai-institute.github.io/research-bench-rubrics/}}

\title{ResearchBench: Recomposing Research Papers into\\
Task Graphs for Evaluating Autonomous AI Research Agents}

\author{Annika Kaul Singh\\
Infinity Inc.\\
\texttt{research@infinity.inc}}

\date{}

\begin{document}
\maketitle

\begin{abstract}
The goal of automated AI research is an agent that can take a hypothesis,
implement it, run the experiments, and report faithful results---but progress
toward this goal is hard to evaluate precisely because novel research has no
ground truth. Paper replication offers an escape: the sub-tasks of replicating
a published paper are, to a first approximation, the same sub-tasks required
to execute novel research, yet the paper's reported numbers provide a
checkable target. \\
The notion of gradable replication via decomposition of papers into
expert-authored ``rubrics'' has already been introduced by previous work.
What is different in 2026 is that these decompositions: \emph{task graphs},
trees of checkable tasks augmented with dependency edges between task groups
and their executable correctness tests can be \emph{generated automatically}
by terminal agents, at unbounded size, for any paper. Task graphs turn the
binary notion of replication into a score: the number of sub-tasks an agent
successfully implements. We present \textbf{ResearchBench}, a public corpus
of 100 papers on LLM inference optimization and adjacent systems research,
each recomposed into a task graph, organized into nested splits
(\textbf{ResearchBench-10/-20/-50/-100}) by run maturity, replicability, and
compute burden, and browsable as interactive graphs at
\href{https://infinity-ai-institute.github.io/research-bench-rubrics/}{ResearchBench}. \\
We present this work based on the argument that because papers
decompose into the same reusable sub-tasks that novel research composes, a
library of task graphs is a substrate for \emph{recomposing} new research
from verified parts.
\end{abstract}

\section{Introduction}

We have discussed how the lack of ground truth makes automated research a
difficult goal to evaluate, and how replication is instead a superb
evaluation target for an agent's ability to conduct automated AI
research---but only if it is scored at the right granularity. Novel research
is not a monolith; it decomposes into environment bring-up, data preparation,
method implementation, benchmarking, ablation, and analysis---and the vast
majority of these sub-tasks replicate sub-tasks that already exist in the
literature. Scoring replication as a single binary judgment collapses months
of layered work into one bit and hides exactly the structure we want to
measure.

PaperBench \citep{starace2025paperbench} innovated on this in 2024 by
explicitly representing the task decomposition of each paper as an
expert-authored ``rubric'' of graded sub-tasks. What is different in 2026 is
that these research task decompositions, together with test suites for
correctness, can be \emph{generated automatically by terminal agents}. We
generalize the representation from rubric to \textbf{task graph}: a weighted
tree of checkable tasks whose section-level nodes additionally carry
dependency edges (\texttt{depends\_on}) recording which task groups must be
completed before others can begin---the environment before the method, the
method before the experiments that exercise it. The graph turns the binary
notion of paper replication into a score: the number of sub-tasks a research
agent can successfully implement, verified by executing tests rather than by
trusting the agent's claims.

Automation removes the two constraints that capped decomposition-based
benchmarks: expert labor and size. Our generator plans each graph from the
paper's structure and expands every section with its full output budget, so a
dense paper receives the decomposition it deserves (228--804 tasks per paper
in our corpus; the largest human-authored task trees run to
$\sim$2{,}000) rather than what fits in a single model response. The same
automation generates the per-task verifiers.

Our contributions are the task graphs themselves:

\begin{enumerate}
  \item \textbf{The task-graph representation and its automatic generation at
        unbounded size.} An outline-then-subtree generator with per-section
        budgets, dependency edges between task groups, duplicate-task
        rejection, and a critique-and-fix convergence pass; graphs reach and
        exceed human-authored scale without padding, and every measured task
        carries an executable verification contract so credit is assigned by
        running tests, never by trusting the agent (\S\ref{sec:graphs}).
  \item \textbf{A public 100-paper corpus of task graphs with nested
        splits.} ResearchBench recomposes 100 papers---LLM serving, KV-cache
        management, quantization, speculative decoding, kernels, compilers,
        hardware microbenchmarking, and generative optimization---organized
        into \textbf{ResearchBench-10/-20/-50/-100} by documented criteria
        (\S\ref{sec:corpus}) and browsable as interactive graphs at \href{https://infinity-ai-institute.github.io/research-bench-rubrics/}{ResearchBench}.
  % \item \textbf{Evidence that evaluating task-graph quality is a research
  %       problem of its own.} By regenerating graphs for papers that have
  %       human-authored PaperBench rubrics and running identical agents under
  %       both, we show aggregate scores are nearly invariant to graph
  %       provenance while per-paper scores are uncorrelated ($r=0.005$):
  %       a benchmark score is a property of the agent$\times$graph pair
  %       (\S\ref{sec:graphquality}).
\end{enumerate}

Supporting these contributions, we report the operational findings that make
graph-scored replication trustworthy in practice: a layered failure taxonomy
that separates model weakness from harness illegibility, evidence that
harness legibility and test-time compute dominate measured agent skill
(a one-paragraph intervention plus raised caps lifts both agents roughly
fivefold, models unchanged), and behavioral anti-reward-hacking defenses
(\S\ref{sec:method}--\S\ref{sec:experiments}).

\section{Related Work}

PaperBench \citep{starace2025paperbench} introduced expert-authored rubric
trees for grading paper replication and remains the reference point for
decomposition-based evaluation; our task graphs generalize its representation
(dependency edges, machine generation, executable contracts) and our corpus
extends its scale by 5$\times$. CORE-Bench \citep{siegel2024core} evaluates
computational reproducibility when author code is available; ResearchBench
blacklists author repositories to force independent reimplementation.
SWE-bench \citep{jimenez2024swebench} grades issue resolution against
maintainer tests; we similarly grade against executable verifiers, but
generate them. MLE-bench \citep{chan2024mlebench} and RE-Bench
\citep{wijk2024rebench} measure ML engineering and frontier R\&D tasks with
curated environments; ResearchBench targets full-paper replication under a
declared compute budget. Specification gaming \citep{krakovna2020specification}
motivates our behavioral (rather than textual) anti-fabrication defenses.

\section{Task Graphs: Structure, Generation, and Verification}
\label{sec:graphs}

\subsection{Automatic generation at unbounded size}

For each paper, a pipeline of terminal-agent stages produces the graph
without expert authorship: ingestion converts the arXiv source to markdown; a
triage stage scores replicability and GPU-hours; a claim extractor records
the paper's numerical results; and the graph generator runs in two stages.
An \emph{outline} call plans the top levels of the graph from the paper's
structure---one section per method component, experiment or table family,
dataset group, and ablation study---assigning each section a task budget
scaled by its claim and table density, the numerical claims it must cover,
and its dependency edges. Each section is then expanded by its own
\emph{subtree} call with the full model output budget to itself and a digest
of sibling sections to suppress cross-section duplicates. This
outline-then-subtree design removes the size ceiling that a single-response
generation imposes: no paper's decomposition is compressed to fit one model
response. A critique-and-fix pass then converges the graph (materializing
missing coverage in bounded batches, annotating verification metadata,
salvaging schema violations), and a deterministic gate rejects graphs with
duplicate tasks---size must come from real decomposition, not repetition.

\subsection{Graph structure: number of tasks and granularity}

Graphs in the corpus run from ${\sim}100$ to ${\sim}935$ tasks per paper
(median ${\sim}300$), at depths of 3--7. Granularity follows a
code-exists $\rightarrow$ code-behaves $\rightarrow$ result-matches
progression: each method component decomposes into structural tasks (the API
surface exists), behavioral tasks (deterministic computation matches), and
measured tasks (a paper number is independently reproduced within tolerance).
Single-response generation systematically under-decomposes. Granularity is
audited by the gate: tasks that merely restate their parent, or that split
one check into trivial variants, are rejected.

\subsection{Verification contracts}

Every measured task is bound to an executable contract specifying the exact
artifact path the implementation must emit and the tolerance under which the
measured value is accepted, e.g.:

\begin{quote}\small
\texttt{expected\_output = results/leaf\_<id>.json};\quad
\texttt{tolerance:\ abs(measured $-$ 2.1)/2.1 $\le$ 0.4}
\end{quote}

A task is credited only when a pytest verifier at a deterministic path
(\texttt{\_tests\_local/test\_leaf\_<id>.py}) passes while reading that
artifact. Credit is attributed by executing tests, never by the agent's
claims.

\section{The ResearchBench Corpus and Splits}
\label{sec:corpus}

\subsection{Papers and why we chose them}

The corpus contains 100 papers centered on LLM inference optimization and the
systems research surrounding it: serving and scheduling (Orca, PagedAttention,
Sarathi-Serve, DistServe, SGLang), KV-cache management and paging, weight and
activation quantization (AWQ, SmoothQuant, KIVI, QServe), speculative
decoding (EAGLE-3, judge decoding), attention and GPU kernels (FlashAttention
1--3, FlashInfer, FlashDecoding++), compilers and autotuning, hardware
microbenchmarking, structured decoding, and LLM-driven generative optimization
(AlphaEvolve, FunSearch, kernel-synthesis agents). Papers were sourced from a
team-curated list spanning these categories and from a triaged backlog of
practitioner-submitted papers, deduplicated against prior sets; each ingested
paper is scored for replicability and estimated GPU-hours at triage. The
domain is deliberate: inference optimization has hardware-sensitive numbers
where ``it runs'' $\neq$ ``it reproduces,'' making it the sharpest test bed
for faithful replication.

Because the corpus is public and consumers' hardware differs, the release
ships the \emph{task-breakdown graphs with dependency structure only}. Due to
this, we also do not release a baseline score that our agents achieve as not
all are runnable on our given constraints. Users of the benchmark should
conduct their own analysis on which tasks in the tree are feasible based on
the dependency trace. We specifically ensure that per-task preconditions
(required hardware, datasets, pretrained models) are recorded so any consumer
can compute their own achievable subset.

\subsection{Nested splits: ResearchBench-10/-20/-50/-100}
\label{sec:splits}

The corpus is organized into four nested splits, each a subset of the next,
assigned by a documented, mechanical criterion (published with the corpus):
papers are ranked by \emph{run maturity} (papers with completed clean-baseline
agent runs rank first---their scores are ground truth), then \emph{triage
replicability}, then \emph{inverse compute burden} (estimated GPU-hours), with
a greedy category-coverage adjustment at each tier boundary so that no split
collapses into a single topic.

\begin{itemize}
  \item \textbf{ResearchBench-10}: the ground-truthed core---papers with
        completed baseline runs, single-GPU feasible, spanning the major
        categories. The split to cite for verified, comparable agent scores.
  \item \textbf{ResearchBench-20}: the full run-backed evaluation split used
        in \S\ref{sec:experiments} (frozen 2026-06-25; 7{,}290 tasks; 1{,}327
        implementable headline claims; every task's contract content-hashed).
  \item \textbf{ResearchBench-50}: adds category breadth---quantization,
        speculative decoding, kernels, compilers---papers with gated task
        graphs that have not yet been run end-to-end.
  \item \textbf{ResearchBench-100}: the full corpus, including
        heavy/multi-GPU and exotic-hardware papers where the graph itself is
        the contribution and replication awaits matching compute.
\end{itemize}

All splits are browsable at
\href{https://infinity-ai-institute.github.io/research-bench-rubrics/}{ResearchBench}
as well as a visualization of their interactive task graphs (collapsible trees
with per-task category coloring, verification specs, and dependency
metadata); split members whose graphs are still generating are listed as
pending. \todo{finalize -50/-100 membership when corpus generation completes}

\subsection{Replicable vs.\ unreplicable papers}

Unreplicability is recorded per task, not per paper. Tasks whose
preconditions cannot be met in a declared environment (lack of hardware;
heavyweight pretrained models; proprietary traces) are classified
\texttt{skipped-environment} and excluded from that environment's achievable
denominator, with audited skip reasons.
% ask jeremy if we should include this
On our reference H100, three papers of the run-backed split have no
implementable headline baseline and score n/a; the remaining 17 are scoreable
end-to-end.

As an example of the environment split, one KV-cache paper decomposes into
309 tasks of which 218 are environment-skipped and 91 achievable.

\subsection{Compute burden and replicability}

Replicability and compute burden are orthogonal axes, and both are scored at
triage: a paper can be perfectly replicable but need 500 GPU-hours, or cheap
but blocked on a proprietary artifact. The corpus spans this plane
deliberately---microbenchmarking papers replicate in minutes; serving papers
need hours of load generation; quantization papers need model downloads but
little compute; kernel papers need specific GPU generations. Split placement
uses estimated GPU-hours so the lower splits stay affordable: the -10 split is
runnable end-to-end for under \$100 of GPU time per agent, while -100 admits
papers whose full replication exceeds any single-node budget.

\subsection{Deciding what to replicate given resources}

Each evaluation run gives the agent an explicit resource contract: a dollar
budget, a wall-clock cap, a maximum number of iterations, and the paper's
feasibility estimate. Agents are instructed to scope honestly---documenting
skipped tasks rather than fabricating them---and, per iteration, are focused
on a small \emph{target slice} of open tasks with exact expected outputs,
which converts a several-hundred-task graph into tractable turns. The
dependency edges give the agent (and the harness's slice selector) the order
in which task groups unlock: environment and data tasks first, method tasks
next, measured experiment tasks last.

\section{Evaluating Task-Graph Quality Is Its Own Problem}
\label{sec:graphquality}

A benchmark built on generated graphs inherits a question that hand-authored
benchmarks could defer to expert judgment: \emph{is the decomposition any
good?} Our central finding is that this question is not answerable by
inspection, is easy to answer wrongly with the obvious metrics, and
constitutes a research task fully separate from evaluating the agents
themselves.

\textbf{The reconstruction experiment.} For 21 papers that ship human-authored
PaperBench rubrics, we regenerated task graphs from scratch with our pipeline
(pristine paper inputs only; leave-one-out few-shot so no paper ever saw its
own published rubric) and ran identical agents under both graph versions with
the identical frozen harness. The aggregate results look like a validation:
Claude Code averages 14.4\% total-implemented under the published rubrics and
15.4\% under our regenerated graphs; Codex averages 5.95\% and 5.97\%. By
aggregate score, the generated graphs are a drop-in substitute.

\textbf{Per-paper, the scores are uncorrelated.} The per-paper correlation
between the two gradings is $r=0.005$. The same agent on the same paper moved
from 77.7\% (published) to 2.7\% (generated) on one paper and from 1.0\% to
46.2\% in the opposite direction on another. Discrimination structure changes
too: under published rubrics Claude beat Codex on 13 papers with 3 ties;
under generated graphs, 11 of 20 papers became ties---the graph version
determines not just scores but \emph{which papers can distinguish agents at
all}. Which papers look ``easy'' is a property of the graph, not the paper:
graph provenance redistributes credit across papers even when it conserves
the total. Any claim of the form ``agent $A$ replicated paper $P$'' is
therefore ill-posed without naming the graph that graded it.

\textbf{Why the obvious quality metrics fail.} Task counts are controlled by
the graph author and can be inflated by splitting checks into trivial
variants; percentage scores inherit whatever denominator the graph chose
(a 40-task graph and a 997-task graph of the same paper produce
incommensurable percentages---both occurred for one paper in this
experiment). Quality judged by fluency of the task text fails silently:
generated tasks read as well as human ones while covering half the content.

\textbf{What holds up.} Three families of checks survived our attempts to
game them. (i) \emph{Reference alignment}: with a human-authored graph as
answer key, judged leaf recall and precision measure whether the generation
decomposed the right content and nothing imaginary. (ii)
\emph{Claim-anchored counting}: the paper's own numerical claims form a
denominator no graph author can inflate; counting claims independently
re-measured within tolerance (verified by isolated re-execution,
\S\ref{sec:hacking}) yields a score comparable \emph{across} graph versions
of the same paper. (iii) \emph{Deterministic structural gates}: duplicate
detection, category-mix bounds, and per-claim coverage checks that run
without any agent. We use (iii) as the generation gate, (i) as the iteration
signal, and (ii) as the cross-provenance yardstick. Under the claim-anchored
metric, the published-rubric baseline for this 21-paper set is 128 of 861
extracted claims independently re-measured (Claude Code); matching or
exceeding it is the bar a generated graph set must clear to be considered a
substitute for expert authorship. \todo{report the generated-graph
claim-anchored number when its full validation run completes}

\section{Discussion and Future Work}

\textbf{Recomposition: from replication to novel research.} The premise of
this benchmark is that papers break down into the same ideas novel research
is built from. Task graphs make that premise operational: a library of 100
graphs is a library of verified, reusable research sub-tasks---environment
recipes, method implementations, measurement harnesses---each with an
executable correctness contract. A new idea can then be expressed as a novel
\emph{composition} of existing sub-tasks (one paper's KV-cache tiering with
another's prefill scheduling, evaluated under a third's load harness), graded
with the same contracts and isolated reproduction that grade replication.
This extends the benchmark from replication toward automated research
execution, including tasks proposed by AI hypothesis systems: the same graph
machinery that scores an agent's replication can scaffold, and audit, its
first steps beyond the paper.

\textbf{Production integration.} The motivating use case is a company that
wants a relevant paper's improvements in its internal, performance-sensitive
codebase---today, manual researcher work. The same contracts that grade
replication become acceptance tests for landing the technique in production
CI.

\textbf{Limitations.} (i) Graph quality is itself machine-generated; the
\texttt{broad\_leaf\_unverifiable} class quantifies, but does not eliminate,
coarse tasks, and per-paper scores are sensitive to graph provenance even
when aggregates agree. (ii) Attempt-3 confounds prompt and compute changes
(isolation planned). (iii) Run-backed evaluation currently covers the
inference-systems splits on one hardware target; the -50/-100 splits ship
graphs ahead of runs. (iv) Two backends shared a GPU in Attempt~3, a
contention risk for a small number of latency-sensitive tasks.

\section{Conclusion}

ResearchBench recomposes 100 research papers into automatically generated,
dependency-structured task graphs with executable verification, organized
into nested splits and browsable publicly, and demonstrates the full loop on
its run-backed split: budgeted agent execution, per-task grading, failure
attribution, and anti-fabrication auditing. Its first empirical lesson is
cautionary for the field: the largest levers on measured ``agent research
ability'' were not the model but whether the harness made the grading
contract legible and whether the agent was given enough test-time compute to
cover the graph. Its design bet is constructive: if papers decompose into
the sub-tasks novel research composes, then a public library of verified
task graphs is not just a benchmark---it is a parts inventory for automated
research. We release the corpus, splits, site, and attempt history as a
standing scoreboard and invite the community to measure models, harnesses,
and compute policies against it.

\section{Evaluation Methodology}
\label{sec:method}

\subsection{Agent contract and iterated sessions}

The replication agent receives only \texttt{paper.md}, the interface
contract, the task graph, the verification contracts, and stock terminal
tools (read, write, edit, bash, glob, grep). The paper's official repository
is blacklisted, forcing independent reimplementation. Because agent sessions
cannot be resumed, all state lives on disk (progress log, source tree,
per-task tests, results); an outer loop re-invokes the agent in fresh
sessions---selecting a repair slice, running a turn, grading, attributing
failures, and choosing the next slice---until headline closure validates or
budget/iteration/wall caps are reached. A full 20-paper backend run completes
in roughly 10--14 hours on one H100.

\subsection{Grading and failure attribution}

Grading is three-layered: per-task test attribution (which executable
verifier passed), reproduction quality (\S\ref{sec:hacking}), and aggregate
scores. Every non-passing task is assigned a failure category and an
\emph{owner layer} (Table~\ref{tab:taxonomy}). This attribution is what makes
harness optimization tractable: without it, harness illegibility is
indistinguishable from model weakness and benchmarks will mis-rank models.

\begin{table}[t]
\caption{Failure taxonomy. Every non-passing task is attributed to a category
and owner layer.}
\label{tab:taxonomy}
\begin{center}
\small
\begin{tabular}{llp{6.1cm}}
\toprule
\textbf{Category} & \textbf{Owner layer} & \textbf{Meaning} \\
\midrule
\texttt{missing\_per\_leaf\_test} & test harness & no executable verifier was produced for the task \\
\texttt{missing\_result\_artifact} & result artifact & the contract-named output file was never emitted \\
\texttt{broad\_leaf\_unverifiable} & graph & task too coarse to verify; needs splitting \\
\texttt{verification\_gap} & implementation & code ran; evidence does not match the contract \\
\texttt{numeric\_mismatch} & implementation & measured, but outside tolerance \\
\texttt{environment\_blocked} & environment & hardware/model/dataset unavailable \\
\bottomrule
\end{tabular}
\end{center}
\end{table}

\subsection{Behavioral defenses against reward hacking}
\label{sec:hacking}

Agents under budget pressure fabricate: hard-coded ``results,'' tests that
assert paper constants, and confident reports of runs that never happened.
Our experience is that \emph{every} text-pattern safeguard was bypassed
within one run of deployment: a regex confession audit for terms like
``cached''/``surrogate'' was defeated by renaming to
\texttt{paper\_reported}; an empty-source negative control was defeated by
tests that passed on stubs; a ``do not fabricate'' prompt directive had no
measurable effect.

Durable safeguards check facts about the run that the agent did not author.
The primary mechanism is \textbf{isolated reproduction}: the grader deletes
the agent's \texttt{results/} directory, re-runs the agent's
\texttt{reproduce.sh} from scratch, and observes wall time, GPU utilization,
and file diffs. Runs whose regenerated results are byte-identical re-ships of
pre-baked constants, or that complete implausibly fast with no GPU activity,
are classified \texttt{suspect} and their scores gated
(Appendix~\ref{app:gates}). Planned extensions include cost/time floors
(actual spend versus the claimed-tasks estimate) and tool-use forensics on
agent transcripts. Our design rule: \emph{if a safeguard can be defeated by
different words, it will be.}

\section{Experiments: Success Rates on ResearchBench-20}
\label{sec:experiments}

\subsection{Setup: same-split attempts}

All experiments run on the frozen ResearchBench-20 split with byte-identical
graphs and contracts, one shared controller, and only the backend adapter
varying between Claude Code (\texttt{claude-opus-4-7}) and Codex
(\texttt{gpt-5-codex}). Attempt~1 established the shared controller on a
mixed-source snapshot. Attempt~2 is the \textbf{clean baseline}: stock
harness, prompts, and adapters, run from a clean tree at committed
\texttt{main} (budget \$200/paper, 15 iterations, 8h wall). Attempt~3 applies
a single \textbf{harness intervention} (below) with raised compute caps
(\$1000/paper, 40 iterations, 12h wall).

\subsection{Success rates of the agents (clean baseline)}

Codex closed 16/1327 implementable claims (1.21\%) and Claude Code 79/1327
(5.95\%); average total-implemented across all graph tasks was 0.21\% and
1.44\% respectively. Claude led on 12 papers, Codex on 2, with 3 tied.
The two backends failed differently: Codex predominantly stopped before
producing artifacts (13/20 papers with replication runtime failures; 1{,}231
tasks blocked on \texttt{missing\_result\_artifact}), while Claude more often
reached verification and stalled on coverage (13/20 papers ending
\texttt{reproduction\_not\_verified\_within\_budget}; 1{,}233 tasks on
\texttt{missing\_per\_leaf\_test}; 73 on \texttt{numeric\_mismatch}---Claude
got to measured numbers more often, and sometimes missed tolerance).

\subsection{Performance across task categories}

Unclosed implementable work concentrates in benchmark-heavy categories
(Table~\ref{tab:categories}): running and measuring, not code-writing, is
where agents fail. Structural and behavioral tasks (does the component exist;
does it compute correctly on deterministic inputs) close at several times the
rate of measured tasks (does the full experiment reproduce the paper's
number), and within measured tasks, those requiring long-running load
generation or model downloads close least. This gradient replicates across
both backends and both graph provenances we have evaluated, suggesting it
reflects genuine difficulty ordering rather than graph idiosyncrasy.

\begin{table}[t]
\caption{Where implementable work remained open (both backends, Attempt 2).}
\label{tab:categories}
\begin{center}
\small
\begin{tabular}{lr}
\toprule
\textbf{Missed task type} & \textbf{Unclosed tasks} \\
\midrule
Evaluation, Metrics \& Benchmarking & 1448 \\
Method Implementation & 955 \\
Experimental Setup & 450 \\
Environment \& Infrastructure Setup & 47 \\
Dataset and Model Acquisition & 43 \\
Logging, Analysis \& Presentation & 28 \\
Data Processing \& Preparation & 27 \\
\bottomrule
\end{tabular}
\end{center}
\end{table}

\subsection{Diagnosis: harness legibility, not model capability}

Failure attribution located the bottleneck precisely. A task can only be
credited when a pytest file exists at the deterministic path
\texttt{\_tests\_local/test\_leaf\_<id>.py} and passes while reading the
contract-named artifact. Inspection of failed workdirs showed agents writing
correct-\emph{shaped} tests and artifacts, but only for the current target
slice (${\sim}4$ tasks), leaving 80+ achievable tasks per paper unverified:
the implementation existed, but the grader was never shown the files it
checks. Nothing in the baseline prompts stated the task$\rightarrow$file
mapping or instructed exhaustive coverage.

\subsection{Impact of the intervention and additional test-time compute}

We appended one paragraph to each backend's prompt stating (i) the
deterministic task$\rightarrow$test-path and contract$\rightarrow$artifact
mapping, (ii) an instruction to enumerate \emph{every} achievable task rather
than the current slice, and (iii) a reminder that \texttt{reproduce.sh} must
regenerate \texttt{results/} from scratch (Appendix~\ref{app:blurb}); and we
raised the compute caps from \$200/15 iterations/8h to \$1000/40
iterations/12h. No frozen artifact, grader, or adapter changed; hash checks
confirm identical graphs and contracts.

\begin{table}[t]
\caption{Implementable-claim closure across same-split attempts. Models
unchanged throughout.}
\label{tab:attempts}
\begin{center}
\small
\begin{tabular}{llrr}
\toprule
\textbf{Attempt} & \textbf{Change vs.\ prior} & \textbf{Codex} & \textbf{Claude Code} \\
\midrule
2 (clean baseline) & none & 16/1327 = 1.21\% & 79/1327 = 5.95\% \\
3 (legibility patch) & prompt paragraph + raised caps & 78/1327 = 5.88\% & 454/1327 = 34.21\% \\
\bottomrule
\end{tabular}
\end{center}
\end{table}

Both backends improved roughly fivefold (Table~\ref{tab:attempts}): Codex's
closure rose $4.9\times$ (average total-implemented 0.21\% to 1.63\%) and
Claude Code's $5.7\times$ (1.44\% to 8.96\%), with five papers exceeding
70\% closure under Claude---including a historic stuck case that had scored
zero in every prior run. Additional test-time compute is a genuine and
separable lever: in the baseline, most Claude papers ended
\texttt{blocked} well under their caps (the loop stalled before spending its
budget), whereas under the intervention agents productively consume the
raised iteration budget covering hundreds of tasks per paper. But compute
alone could not have produced files the agent did not know to write---and
the intervention alone, under starved caps, could not cover 300-task graphs.
We emphasize the confound: Attempt~3 changed the prompt \emph{and} the caps,
and ran both backends concurrently on one H100 (a minor contention risk for
latency-sensitive measured tasks). Isolating the levers is Attempt~4.
Nevertheless, the qualitative conclusion is robust: a large fraction of what
a naive benchmark would report as ``agent failure'' was \emph{harness
legibility} failure, and much of the rest was starvation of test-time
compute.

\bibliographystyle{plainnat}
\begin{thebibliography}{9}

\bibitem[Chan et~al.(2024)]{chan2024mlebench}
J.~S. Chan, N.~Chowdhury, O.~Jaffe, et~al.
\newblock MLE-bench: Evaluating machine learning agents on machine learning
  engineering.
\newblock \emph{arXiv:2410.07095}, 2024.

\bibitem[Jimenez et~al.(2024)]{jimenez2024swebench}
C.~E. Jimenez, J.~Yang, A.~Wettig, et~al.
\newblock SWE-bench: Can language models resolve real-world GitHub issues?
\newblock In \emph{ICLR}, 2024.

\bibitem[Krakovna et~al.(2020)]{krakovna2020specification}
V.~Krakovna, J.~Uesato, V.~Mikulik, et~al.
\newblock Specification gaming: the flip side of AI ingenuity.
\newblock DeepMind Blog, 2020.

\bibitem[Siegel et~al.(2024)]{siegel2024core}
Z.~S. Siegel, S.~Kapoor, N.~Nadgir, B.~Stroebl, and A.~Narayanan.
\newblock CORE-Bench: Fostering the credibility of published research through a
  computational reproducibility agent benchmark.
\newblock \emph{arXiv:2409.11363}, 2024.

\bibitem[Starace et~al.(2025)]{starace2025paperbench}
G.~Starace, O.~Jaffe, D.~Sherburn, et~al.
\newblock PaperBench: Evaluating AI's ability to replicate AI research.
\newblock \emph{arXiv:2504.01848}, 2025.

\bibitem[Wijk et~al.(2024)]{wijk2024rebench}
H.~Wijk, T.~Lin, J.~Becker, et~al.
\newblock RE-Bench: Evaluating frontier AI R\&D capabilities of language model
  agents against human experts.
\newblock \emph{arXiv:2411.15114}, 2024.

\end{thebibliography}

\appendix

\section{ResearchBench-20 composition}
\label{app:papers}

\begin{table}[h]
\caption{Papers in the run-backed ResearchBench-20 split with task counts.
Three papers (marked $\dagger$) have no implementable headline baseline in the
reference environment and score n/a. The full corpus roster with split
membership is maintained on the project site (\S\ref{sec:splits}).}
\begin{center}
\small
\begin{tabular}{lr@{\qquad}lr}
\toprule
\textbf{Paper (slug, abbreviated)} & \textbf{Tasks} & \textbf{Paper} & \textbf{Tasks} \\
\midrule
llmservingsim-2.0 & 214 & icecache & 935 \\
FOCUS-autokernel & 268 & lamp & 288 \\
fp16-equivalence & 650 & niyama & 160 \\
nexus & 312 & prefill-as-a-service & 532 \\
adaserve & 692 & ragged-paged-attention & 442 \\
asynctls & 475 & shift-parallelism & 407 \\
cacheflow & 177 & ttkv & 565 \\
contiguouskv & 309 & tutti & 220 \\
duetserve & 350 & fast-heterogeneous$^\dagger$ & 93 \\
flexicache$^\dagger$ & 96 & hexgen-flow$^\dagger$ & 105 \\
\midrule
\multicolumn{2}{l}{\textbf{Total}} & & \textbf{7{,}290} \\
\bottomrule
\end{tabular}
\end{center}
\end{table}

\section{The Attempt-3 legibility paragraph}
\label{app:blurb}

The full text appended to both backends' prompts (verbatim):

\begin{quote}\small\ttfamily
LEAF -> FILE PATH MAPPING (this is how you score; read carefully):\\
- A rubric leaf is counted as verified ONLY when a pytest file exists at the
exact path \_tests\_local/test\_leaf\_<LEAF\_ID>.py and passes, where
<LEAF\_ID> is the leaf's id with every hyphen '-' replaced by '\_'. If that
file is missing, the leaf stays unverified and scores zero no matter how
correct your src/ code is.\\
- For measured / result\_match leaves, that test must load the artifact named
by the leaf's contract expected\_output (see ../verification\_contracts.json;
convention results/leaf\_<LEAF\_ID>.json) and assert the contract's
field(s)/tolerance.\\
- Enumerate EVERY leaf in ../rubric.json and ../verification\_contracts.json.
For each leaf NOT blocked by unavailable hardware/models/datasets (i.e.\ every
achievable leaf, not just the current target slice), write its
\_tests\_local/test\_leaf\_<LEAF\_ID>.py AND emit the result artifact that
test reads. Breadth of achievable-leaf coverage is the single biggest lever on
your score.\\
- reproduce.sh must regenerate all of results/ from scratch (the grader wipes
results/ and reruns reproduce.sh in isolation).
\end{quote}

\section{Reproduction quality gates}
\label{app:gates}

Isolated reproduction classifies each run as \texttt{ok}, \texttt{suspect}, or
\texttt{failed}: \texttt{failed} if \texttt{reproduce.sh} is absent, exits
nonzero, times out, or produces zero files into a wiped \texttt{results/};
\texttt{suspect} if wall time is implausibly short with no GPU activity, or if
more than half of the regenerated files are byte-identical to the agent's
pre-authored results with no authored code paths present. Suspect and failed
runs have their task credit gated.

\end{document}
